% ============================================================
%  ACADEMIA ZENIT — GUÍA COMPLETA: NÚMEROS ENTEROS Y RACIONALES
%  Matemática · 1° Medio
% ============================================================

\documentclass[11pt, letterpaper]{article}

% -- Paquetes de Compatibilidad y Corrección --
\usepackage[T1]{fontenc}
\usepackage[utf8]{inputenc}
\usepackage[spanish, es-tabla]{babel}

\usepackage[margin=2cm, top=2.2cm, bottom=2.2cm]{geometry}
\usepackage{amsmath, amssymb, amsthm}
\usepackage{mathtools}
\usepackage{array, booktabs, tabularx}
\usepackage{xcolor}
\usepackage{graphicx}
\usepackage{enumitem}
\usepackage{fancyhdr}
\usepackage{tcolorbox}
\usepackage{tikz}
\usepackage{pgfplots}
\pgfplotsset{compat=1.18}
\usepackage{lastpage}
\usepackage{multicol}
\usepackage{cancel}

% -- Colores Institucionales Academia Zenit --
\definecolor{zenitnavy}{HTML}{1A3A5C}
\definecolor{zenitblue}{HTML}{2E6DA4}
\definecolor{zenitgold}{HTML}{C9820A}
\definecolor{zenitlightgray}{HTML}{F5F5F5}
\definecolor{zenitgreen}{HTML}{1A6B3C}
\definecolor{zenitred}{HTML}{8B1A1A}
\definecolor{blocka}{HTML}{1A3A5C}
\definecolor{blockb}{HTML}{1A6B3C}
\definecolor{blockc}{HTML}{7B2D8B}
\definecolor{teorema}{HTML}{E8F4FD}
\definecolor{ejemplo}{HTML}{FFF8E7}
\definecolor{ejercicio}{HTML}{F0FFF0}
\definecolor{alerta}{HTML}{FFF0F0}

% -- Estilos de tcolorbox --
\tcbuselibrary{skins, breakable}

\tcbset{
  instrbox/.style={
    enhanced, colback=white, colframe=zenitgold,
    left=8pt, right=8pt, top=6pt, bottom=6pt,
    arc=2pt, leftrule=4pt,
  },
  teoremabox/.style={
    enhanced, colback=teorema, colframe=zenitblue,
    left=8pt, right=8pt, top=6pt, bottom=6pt,
    arc=4pt, leftrule=5pt,
    fonttitle=\bfseries\color{zenitblue},
  },
  ejemplobox/.style={
    enhanced, colback=ejemplo, colframe=zenitgold,
    left=8pt, right=8pt, top=6pt, bottom=6pt,
    arc=4pt, leftrule=5pt,
    fonttitle=\bfseries\color{zenitgold},
  },
  ejerciciobox/.style={
    enhanced, colback=ejercicio, colframe=zenitgreen,
    left=8pt, right=8pt, top=6pt, bottom=6pt,
    arc=4pt, leftrule=5pt,
    fonttitle=\bfseries\color{zenitgreen},
  },
  alertabox/.style={
    enhanced, colback=alerta, colframe=zenitred,
    left=8pt, right=8pt, top=6pt, bottom=6pt,
    arc=4pt, leftrule=5pt,
    fonttitle=\bfseries\color{zenitred},
  },
}

% -- Encabezado de bloque --
\newcommand{\bloqueheader}[3]{%
  \vspace{10pt}
  \noindent
  \begin{tikzpicture}[baseline]
    \fill[#2] (0,0) rectangle (2.5cm, 0.7cm);
    \node[white, font=\bfseries] at (1.25cm, 0.35cm) {Bloque #1};
  \end{tikzpicture}%
  \begin{tikzpicture}[baseline]
    \fill[#2!10!white] (0,0) rectangle (\linewidth - 2.6cm, 0.7cm);
    \node[color=#2, font=\bfseries\large, anchor=west] at (0.2cm, 0.35cm) {#3};
  \end{tikzpicture}%
  \vspace{4pt}
}

\newcommand{\subbloque}[1]{%
  \vspace{8pt}
  \noindent\textbf{\color{zenitnavy}\large #1}
  \vspace{3pt}
  \hrule height 1pt
  \vspace{5pt}
}

% -- Configuración de Encabezado --
\pagestyle{fancy}
\fancyhf{}
\fancyhead[L]{\small \textbf{Academia Zenit} $\cdot$ Números Enteros y Racionales}
\fancyhead[R]{\small Página \thepage\ de \pageref{LastPage}}
\fancyfoot[L]{\small \textit{1° Medio · Matemática}}
\fancyfoot[R]{\small Prohibida su reproducción total o parcial.}

% ============================================================
\begin{document}
% ============================================================

\begin{center}
  {\Huge\bfseries\color{zenitnavy} ACADEMIA ZENIT}\\[4pt]
  {\LARGE\bfseries\color{zenitblue} GUÍA COMPLETA: NÚMEROS ENTEROS Y RACIONALES}\\[4pt]
  {\large\color{zenitgold} Matemática · 1° Medio}\\[6pt]
  \rule{\linewidth}{2pt}
\end{center}

\begin{tcolorbox}[instrbox]
  \textbf{¿Cómo usar esta guía?} Encontrarás primero la materia con ejemplos resueltos paso a paso. Luego, ejercicios para practicar. Al final, una mezcla de todos los temas. Resuelve cada ejercicio antes de ver la respuesta. ¡El secreto está en la práctica!
\end{tcolorbox}

\vspace{6pt}

% ============================================================
% --- BLOQUE A: NÚMEROS ENTEROS — FUNDAMENTOS ---
% ============================================================
\bloqueheader{A}{blocka}{Números Enteros: Fundamentos y Operaciones Básicas}

\subbloque{1. ¿Qué son los números enteros?}

\begin{tcolorbox}[teoremabox, title={Definición}]
El conjunto de los \textbf{números enteros} es $\mathbb{Z} = \{\ldots, -3,\,-2,\,-1,\,0,\,1,\,2,\,3,\ldots\}$.
Incluye los números \textbf{positivos}, el \textbf{cero} y los números \textbf{negativos}.
\end{tcolorbox}

En la recta numérica, los números enteros se ubican simétricamente respecto al cero:

\begin{center}
\begin{tikzpicture}
  \draw[-latex, thick] (-5.5,0) -- (5.5,0);
  \foreach \x in {-5,-4,-3,-2,-1,0,1,2,3,4,5} {
    \draw (\x, 0.1) -- (\x, -0.1);
    \node[below] at (\x, -0.1) {\small $\x$};
  }
  \foreach \x in {-4,-3,-2,-1} \fill[zenitred] (\x,0) circle (2.5pt);
  \fill[zenitnavy] (0,0) circle (2.5pt);
  \foreach \x in {1,2,3,4} \fill[zenitgreen] (\x,0) circle (2.5pt);
\end{tikzpicture}
\end{center}

\subbloque{2. Suma y Resta de Enteros}

\begin{tcolorbox}[teoremabox, title={Reglas fundamentales}]
\begin{enumerate}[leftmargin=*, itemsep=3pt]
  \item \textbf{Mismos signos:} se suman los valores absolutos y se conserva el signo.
  \item \textbf{Signos distintos:} se restan los valores absolutos y se toma el signo del mayor.
  \item \textbf{Restar un número} equivale a \textbf{sumar su opuesto}: $\;a - (-b) = a + b$
\end{enumerate}
\end{tcolorbox}

\begin{tcolorbox}[alertabox, title={¡Ojo con los signos!}]
$-(-3) = +3$ \quad y \quad $-(+3) = -3$ \quad. Cuando hay dos signos juntos: \textbf{iguales} $\to$ positivo; \textbf{distintos} $\to$ negativo.
\end{tcolorbox}

\begin{tcolorbox}[ejemplobox, title={Ejemplos resueltos}]
\textbf{a)} $(-8) + (-5) = -(8+5) = -13$ \hfill {\small (mismos signos, ambos negativos)}\\[4pt]
\textbf{b)} $(-7) + (+4) = -(7-4) = -3$ \hfill {\small (distintos signos, mayor valor absoluto es 7)}\\[4pt]
\textbf{c)} $(-3) - (+6) = -3 + (-6) = -9$ \hfill {\small (restar positivo = sumar negativo)}\\[4pt]
\textbf{d)} $(-3) - (-6) = -3 + 6 = +3$ \hfill {\small (restar negativo = sumar positivo)}\\[4pt]
\textbf{e)} $5 - 9 + 3 - (-2) = 5 - 9 + 3 + 2 = 1$ \hfill {\small (varias operaciones)}
\end{tcolorbox}

\begin{tcolorbox}[ejerciciobox, title={Ejercicios — Nivel Básico (Suma y Resta)}]
Calcula el resultado de cada operación:
\begin{multicols}{2}
\begin{enumerate}[label=\textbf{\arabic*.}, itemsep=6pt]
  \item $(-4) + (-9) = $
  \item $(-12) + 5 = $
  \item $8 + (-15) = $
  \item $(-6) - (-6) = $
  \item $(-3) - (+7) = $
  \item $0 - (-11) = $
  \item $14 + (-14) = $
  \item $(-20) - (-8) = $
  \item $(-5) + 3 + (-7) = $
  \item $12 - 7 + (-4) + 3 = $
\end{enumerate}
\end{multicols}
\vspace{2pt}
%\small\textit{respuestas: 1) $-13$ \; 2) $-7$ \; 3) $-7$ \; 4) $0$ \; 5) $-10$ \; 6) $11$ \; 7) $0$ \; 8) $-12$ \; 9) $-9$ \; 10) $4$}
\end{tcolorbox}

\subbloque{3. Multiplicación y División de Enteros}

\begin{tcolorbox}[teoremabox, title={Ley de los signos}]
\begin{center}
\renewcommand{\arraystretch}{1.3}
\begin{tabular}{c|c|c}
\textbf{Factor 1} & \textbf{Factor 2} & \textbf{Resultado} \\
\hline
$(+)$ & $(+)$ & $(+)$ \\
$(+)$ & $(-)$ & $(-)$ \\
$(-)$ & $(+)$ & $(-)$ \\
$(-)$ & $(-)$ & $(+)$ \\
\end{tabular}
\end{center}
\textbf{Regla práctica:} cantidad \textbf{par} de signos negativos $\to$ resultado positivo; cantidad \textbf{impar} $\to$ resultado negativo.
\end{tcolorbox}

\begin{tcolorbox}[ejemplobox, title={Ejemplos resueltos}]
\textbf{a)} $(-4) \cdot (-3) = +12$ \hfill {\small (dos negativos $\to$ positivo)}\\[4pt]
\textbf{b)} $(-5) \cdot (+7) = -35$ \hfill {\small (un negativo $\to$ negativo)}\\[4pt]
\textbf{c)} $(-2) \cdot (-3) \cdot (-4) = +6 \cdot (-4) = -24$ \hfill {\small (tres negativos $\to$ negativo)}\\[4pt]
\textbf{d)} $\dfrac{-36}{-9} = +4$ \hfill {\small (mismos signos $\to$ positivo)}\\[6pt]
\textbf{e)} $\dfrac{-48}{+6} = -8$ \hfill {\small (signos distintos $\to$ negativo)}
\end{tcolorbox}

\begin{tcolorbox}[ejerciciobox, title={Ejercicios — Nivel Básico (Multiplicación y División)}]
\begin{multicols}{2}
\begin{enumerate}[label=\textbf{\arabic*.}, itemsep=6pt]
  \item $(-6) \cdot (-7) = $
  \item $(-8) \cdot (+5) = $
  \item $(+9) \cdot (-4) = $
  \item $(-3) \cdot (-3) \cdot (-3) = $
  \item $\dfrac{-56}{+8} = $
  \item $\dfrac{-72}{-9} = $
  \item $\dfrac{+45}{-5} = $
  \item $(-2) \cdot (-5) \cdot (+4) = $
  \item $(-1) \cdot (-1) \cdot (-1) \cdot (-1) = $
  \item $\dfrac{(-3) \cdot (-8)}{-6} = $
\end{enumerate}
\end{multicols}
\vspace{2pt}
%\small\textit{respuestas: 1) $42$ \; 2) $-40$ \; 3) $-36$ \; 4) $-27$ \; 5) $-7$ \; 6) $8$ \; 7) $-9$ \; 8) $40$ \; 9) $1$ \; 10) $-4$}
\end{tcolorbox}

% ============================================================
% --- BLOQUE A (CONT): POTENCIAS CON ENTEROS ---
% ============================================================
\bloqueheader{A}{blocka}{Potencias con Números Enteros}

\subbloque{4. Potencias de base entera}

\begin{tcolorbox}[teoremabox, title={Propiedad clave}]
\[
  (-a)^n = \begin{cases} +a^n & \text{si } n \text{ es par} \\ -a^n & \text{si } n \text{ es impar} \end{cases}
\]
\textbf{Atención:} $-a^n \neq (-a)^n$. El signo afuera no eleva; el signo adentro del paréntesis sí.
\end{tcolorbox}

\begin{tcolorbox}[alertabox, title={Diferencia crucial}]
$\mathbf{-3^2 = -(3^2) = -9}$ \hspace{1cm} pero \hspace{1cm} $\mathbf{(-3)^2 = (-3)\cdot(-3) = +9}$
\end{tcolorbox}

\begin{tcolorbox}[ejemplobox, title={Ejemplos resueltos}]
\textbf{a)} $(-2)^4 = (-2)\cdot(-2)\cdot(-2)\cdot(-2) = 4 \cdot 4 = +16$\\[4pt]
\textbf{b)} $(-2)^5 = (-2)^4 \cdot (-2) = 16 \cdot (-2) = -32$\\[4pt]
\textbf{c)} $-5^2 = -(5 \cdot 5) = -25$\\[4pt]
\textbf{d)} $(-5)^2 = +25$\\[4pt]
\textbf{e)} $(-1)^{100} = +1$ \hfill {\small (exponente par)}\\[4pt]
\textbf{f)} $(-1)^{101} = -1$ \hfill {\small (exponente impar)}
\end{tcolorbox}

\begin{tcolorbox}[ejerciciobox, title={Ejercicios — Potencias}]
\begin{multicols}{2}
\begin{enumerate}[label=\textbf{\arabic*.}, itemsep=6pt]
  \item $(-3)^3 = $
  \item $(-4)^2 = $
  \item $-4^2 = $
  \item $(-2)^6 = $
  \item $(-1)^{55} = $
  \item $(-10)^3 = $
  \item $(-5)^4 = $
  \item $-2^4 = $
  \item $(-3)^2 - (-2)^3 = $
  \item $(-1)^{100} + (-1)^{101} = $
\end{enumerate}
\end{multicols}
\vspace{2pt}
%\small\textit{respuestas: 1) $-27$ \; 2) $16$ \; 3) $-16$ \; 4) $64$ \; 5) $-1$ \; 6) $-1000$ \; 7) $625$ \; 8) $-16$ \; 9) $17$ \; 10) $0$}
\end{tcolorbox}

% ============================================================
% --- BLOQUE A: JERARQUÍA DE OPERACIONES ---
% ============================================================
\bloqueheader{A}{blocka}{Jerarquía de Operaciones y Símbolos de Agrupación}

\subbloque{5. Orden correcto para operar}

\begin{tcolorbox}[teoremabox, title={Jerarquía de operaciones (de mayor a menor prioridad)}]
\begin{enumerate}[leftmargin=*, itemsep=4pt]
  \item \textbf{Símbolos de agrupación:} primero el más interno $( ) \to [ ] \to \{ \}$
  \item \textbf{Potencias y raíces}
  \item \textbf{Multiplicaciones y divisiones} (de izquierda a derecha)
  \item \textbf{Sumas y restas} (de izquierda a derecha)
\end{enumerate}
\end{tcolorbox}

\begin{tcolorbox}[ejemplobox, title={Ejemplos resueltos — un símbolo de agrupación}]
\textbf{a)} $3 + 2 \cdot (-4) = 3 + (-8) = -5$ \hfill {\small (primero la multiplicación)}\\[4pt]
\textbf{b)} $(3 + 2) \cdot (-4) = 5 \cdot (-4) = -20$ \hfill {\small (primero el paréntesis)}\\[4pt]
\textbf{c)} $(-3)^2 - 2 \cdot 5 = 9 - 10 = -1$ \hfill {\small (potencia y multiplicación antes de resta)}\\[4pt]
\textbf{d)} $\dfrac{(-4)^2 + 2}{3 - 5} = \dfrac{16 + 2}{-2} = \dfrac{18}{-2} = -9$
\end{tcolorbox}

\begin{tcolorbox}[ejemplobox, title={Ejemplos resueltos — varios símbolos de agrupación}]
Resolver: $\{3 - [2 \cdot (-1 + 4)]^2 + 5\} \cdot (-2)$

\medskip
\textbf{Paso 1:} El paréntesis más interno: $(-1 + 4) = 3$

\textbf{Paso 2:} El corchete con la potencia: $[2 \cdot 3]^2 = [6]^2 = 36$

\textbf{Paso 3:} La llave: $\{3 - 36 + 5\} = \{-28\}$

\textbf{Paso 4:} Multiplicar: $(-28) \cdot (-2) = \mathbf{56}$
\end{tcolorbox}

\begin{tcolorbox}[ejerciciobox, title={Ejercicios — Nivel Intermedio (Jerarquía)}]
\begin{enumerate}[label=\textbf{\arabic*.}, itemsep=8pt]
  \item $5 - 3 \cdot (-2) + 4 =$
  \item $(-2)^3 + 4 \cdot (-1) - 5 =$
  \item $(8 - 11) \cdot (-3) - 2^2 =$
  \item $\dfrac{-12 + 6 \cdot (-1)}{(-3)^2 - 9 + 3} =$
  \item $[(-4) + 7]^2 - 5 \cdot (-2) + 3 =$
  \item $-3 \cdot [(5 - 8)^2 - 1] + 6 =$
  \item $\{[(-2)^3 + 4] \cdot (-3)\} - (-1)^5 =$
  \item $2 \cdot \{-[3 - (4 - 7)] + 2^3\} - 5 =$
\end{enumerate}
\vspace{2pt}
%\small\textit{respuestas: 1) $15$ \; 2) $-17$ \; 3) $5$ \; 4) $-6$ \; 5) $22$ \; 6) $-18$ \; 7) $13$ \; 8) $19$}
\end{tcolorbox}

\begin{tcolorbox}[ejerciciobox, title={Ejercicios — Nivel Avanzado (Múltiples Agrupaciones con Enteros)}]
Resuelve completamente, mostrando cada paso:
\begin{enumerate}[label=\textbf{\arabic*.}, itemsep=10pt]
  \item $\{(-3)^3 - [(-2)^4 - 5 \cdot (-1+3)^2] + 8\} \div (-5) =$
  \item $-\{2 \cdot [(-1)^5 + 3 \cdot (-2 + 5)]^2 - 4^2\} + (-3)^2 =$
  \item $\dfrac{\{3 - [(-2)^3 + 5] \cdot (-1)\}^2}{(-7 + 3)^2 - 4^2 + 8 \cdot (-2)} =$
  \item $\big\{[(-4)^2 - 3 \cdot (-1-2)^2] \cdot (-2)\big\} - \{(-2)^5 + [3 - (-1)^3]^2\} =$
  \item $\dfrac{(-3)^4 - [(- 2)^3 \cdot(-1)^7 + 5]^2}{(-1 + 2 \cdot 3)^2 - 4 \cdot 5} =$
  \item $-2 \cdot \left\{-3 + \left[(-2)^2 - \dfrac{(-4)^2 - 8}{(-1)^5 + 3}\right]^3 \right\} + (-5)^2 =$
\end{enumerate}
\vspace{2pt}
%\small\textit{respuestas: 1) $7$ \; 2) $-$ {\normalsize calculalo!} (ver abajo)} \\[4pt]
\begin{tcolorbox}[colback=white, colframe=zenitnavy, arc=2pt, left=4pt, right=4pt, top=3pt, bottom=3pt]
\textbf{Solución detallada del \#1:}
\begin{align*}
&\{(-3)^3 - [(-2)^4 - 5 \cdot (-1+3)^2] + 8\} \div (-5) \\
&= \{-27 - [16 - 5 \cdot (2)^2] + 8\} \div (-5) \\
&= \{-27 - [16 - 5 \cdot 4] + 8\} \div (-5) \\
&= \{-27 - [16 - 20] + 8\} \div (-5) \\
&= \{-27 - (-4) + 8\} \div (-5) \\
&= \{-27 + 4 + 8\} \div (-5) \\
&= (-15) \div (-5) = \mathbf{3}
\end{align*}
\end{tcolorbox}
\end{tcolorbox}

\newpage

% ============================================================
% --- BLOQUE B: NÚMEROS RACIONALES ---
% ============================================================
\bloqueheader{B}{blockb}{Números Racionales: Fundamentos}

\subbloque{6. ¿Qué son los números racionales?}

\begin{tcolorbox}[teoremabox, title={Definición}]
Un \textbf{número racional} es todo número que puede escribirse como $\dfrac{p}{q}$ con $p, q \in \mathbb{Z}$ y $q \neq 0$.
El conjunto se denota $\mathbb{Q}$. Incluye fracciones, números enteros, decimales finitos y periódicos.
\end{tcolorbox}

\begin{tcolorbox}[alertabox, title={¡No dividir por cero!}]
$\dfrac{5}{0}$ \textbf{no existe}. La división por cero no está definida en ningún contexto.
\end{tcolorbox}

\subbloque{7. Fracciones equivalentes y simplificación}

\begin{tcolorbox}[teoremabox, title={Fracción equivalente}]
$\dfrac{a}{b} = \dfrac{a \cdot k}{b \cdot k}$ para cualquier $k \neq 0$. \quad
Una fracción está en su \textbf{mínima expresión} cuando $\text{mcd}(|p|, |q|) = 1$.
\end{tcolorbox}

\begin{tcolorbox}[ejemplobox, title={Ejemplos resueltos}]
\textbf{a)} Simplificar $\dfrac{24}{36}$: \quad $\text{mcd}(24,36)=12$ \quad $\Rightarrow$ \quad $\dfrac{24\div12}{36\div12}=\dfrac{2}{3}$\\[8pt]
\textbf{b)} Simplificar $\dfrac{-45}{30}$: \quad $\text{mcd}(45,30)=15$ \quad $\Rightarrow$ \quad $\dfrac{-45\div15}{30\div15}=\dfrac{-3}{2}$\\[8pt]
\textbf{c)} ¿Son equivalentes $\dfrac{3}{4}$ y $\dfrac{15}{20}$? \quad $\dfrac{15\div5}{20\div5}=\dfrac{3}{4}$ \quad \checkmark \quad Sí.
\end{tcolorbox}

\subbloque{8. Suma y Resta de Fracciones}

\begin{tcolorbox}[teoremabox, title={Procedimiento}]
\begin{enumerate}[leftmargin=*, itemsep=3pt]
  \item Si tienen el \textbf{mismo denominador}: $\dfrac{a}{c} \pm \dfrac{b}{c} = \dfrac{a \pm b}{c}$
  \item Si tienen \textbf{distinto denominador}: calcular el \textbf{m.c.m.} y llevar a común denominador.
  \item Simplificar el resultado.
\end{enumerate}
\end{tcolorbox}

\begin{tcolorbox}[ejemplobox, title={Ejemplos resueltos}]
\textbf{a)} $\dfrac{3}{7} + \dfrac{2}{7} = \dfrac{5}{7}$ \hfill {\small (mismo denominador)}\\[8pt]
\textbf{b)} $\dfrac{1}{4} + \dfrac{2}{3}$: \quad $\text{m.c.m.}(4,3) = 12$
\[
  \dfrac{1 \cdot 3}{12} + \dfrac{2 \cdot 4}{12} = \dfrac{3}{12} + \dfrac{8}{12} = \dfrac{11}{12}
\]
\textbf{c)} $\dfrac{5}{6} - \dfrac{3}{8}$: \quad $\text{m.c.m.}(6,8) = 24$
\[
  \dfrac{5 \cdot 4}{24} - \dfrac{3 \cdot 3}{24} = \dfrac{20}{24} - \dfrac{9}{24} = \dfrac{11}{24}
\]
\textbf{d)} $\dfrac{-2}{5} + \dfrac{3}{-4} = \dfrac{-2}{5} - \dfrac{3}{4}$: \quad $\text{m.c.m.}(5,4)=20$
\[
  \dfrac{-8}{20} - \dfrac{15}{20} = \dfrac{-23}{20}
\]
\end{tcolorbox}

\begin{tcolorbox}[ejerciciobox, title={Ejercicios — Suma y Resta de Fracciones}]
\begin{multicols}{2}
\begin{enumerate}[label=\textbf{\arabic*.}, itemsep=8pt]
  \item $\dfrac{5}{9} + \dfrac{1}{9} =$
  \item $\dfrac{3}{4} - \dfrac{1}{4} =$
  \item $\dfrac{1}{3} + \dfrac{1}{4} =$
  \item $\dfrac{5}{6} - \dfrac{1}{4} =$
  \item $\dfrac{-3}{5} + \dfrac{1}{2} =$
  \item $\dfrac{7}{8} - \dfrac{2}{3} =$
  \item $1 - \dfrac{3}{7} =$
  \item $\dfrac{-2}{3} - \dfrac{1}{6} =$
  \item $\dfrac{3}{4} + \dfrac{-5}{8} =$
  \item $\dfrac{5}{12} - \dfrac{3}{8} + \dfrac{1}{6} =$
\end{enumerate}
\end{multicols}
\vspace{2pt}
%\small\textit{respuestas: 1) $\tfrac{2}{3}$ \; 2) $\tfrac{1}{2}$ \; 3) $\tfrac{7}{12}$ \; 4) $\tfrac{7}{12}$ \; 5) $\tfrac{-1}{10}$ \; 6) $\tfrac{5}{24}$ \; 7) $\tfrac{4}{7}$ \; 8) $\tfrac{-5}{6}$ \; 9) $\tfrac{1}{8}$ \; 10) $\tfrac{1}{8}$}
\end{tcolorbox}

\subbloque{9. Multiplicación y División de Fracciones}

\begin{tcolorbox}[teoremabox, title={Reglas}]
\[
  \dfrac{a}{b} \cdot \dfrac{c}{d} = \dfrac{a \cdot c}{b \cdot d} \qquad\qquad
  \dfrac{a}{b} \div \dfrac{c}{d} = \dfrac{a}{b} \cdot \dfrac{d}{c} = \dfrac{a \cdot d}{b \cdot c}
\]
\textbf{Dividir} por una fracción es \textbf{multiplicar} por su \textbf{recíproca} (fracción invertida).
\end{tcolorbox}

\begin{tcolorbox}[ejemplobox, title={Ejemplos resueltos}]
\textbf{a)} $\dfrac{3}{4} \cdot \dfrac{8}{9} = \dfrac{3 \cdot 8}{4 \cdot 9} = \dfrac{24}{36} = \dfrac{2}{3}$\\[8pt]
\textbf{b)} $\dfrac{-5}{6} \cdot \dfrac{3}{10} = \dfrac{-15}{60} = \dfrac{-1}{4}$\\[8pt]
\textbf{c)} $\dfrac{2}{3} \div \dfrac{4}{9} = \dfrac{2}{3} \cdot \dfrac{9}{4} = \dfrac{18}{12} = \dfrac{3}{2}$\\[8pt]
\textbf{d)} $\dfrac{-7}{8} \div \dfrac{-21}{4} = \dfrac{-7}{8} \cdot \dfrac{4}{-21} = \dfrac{28}{168} = \dfrac{1}{6}$
\end{tcolorbox}

\begin{tcolorbox}[ejerciciobox, title={Ejercicios — Multiplicación y División de Fracciones}]
\begin{multicols}{2}
\begin{enumerate}[label=\textbf{\arabic*.}, itemsep=8pt]
  \item $\dfrac{2}{3} \cdot \dfrac{9}{4} =$
  \item $\dfrac{-5}{7} \cdot \dfrac{14}{15} =$
  \item $\dfrac{3}{8} \cdot \dfrac{-4}{9} =$
  \item $\dfrac{6}{5} \div \dfrac{3}{10} =$
  \item $\dfrac{-4}{7} \div \dfrac{8}{21} =$
  \item $\dfrac{-3}{4} \div \left(-\dfrac{9}{8}\right) =$
  \item $\dfrac{5}{6} \cdot \dfrac{3}{5} \cdot \dfrac{2}{9} =$
  \item $\left(\dfrac{-2}{3}\right)^2 =$
  \item $\dfrac{4}{3} \cdot \left(\dfrac{-3}{8}\right) \div \dfrac{1}{2} =$
  \item $\dfrac{-5}{4} \div \dfrac{15}{-8} =$
\end{enumerate}
\end{multicols}
\vspace{2pt}
%\small\textit{respuestas: 1) $\tfrac{3}{2}$ \; 2) $\tfrac{-2}{3}$ \; 3) $\tfrac{-1}{6}$ \; 4) $4$ \; 5) $\tfrac{-3}{2}$ \; 6) $\tfrac{2}{3}$ \; 7) $\tfrac{1}{9}$ \; 8) $\tfrac{4}{9}$ \; 9) $\tfrac{-1}{1}=-1$ \; 10) $\tfrac{2}{3}$}
\end{tcolorbox}

% ============================================================
% --- BLOQUE B (CONT): POTENCIAS CON RACIONALES ---
% ============================================================
\bloqueheader{B}{blockb}{Potencias y Expresiones con Racionales}

\subbloque{10. Potencias de base racional y exponente negativo}

\begin{tcolorbox}[teoremabox, title={Exponente negativo}]
\[
  \left(\dfrac{a}{b}\right)^{-n} = \left(\dfrac{b}{a}\right)^{n} \qquad (a, b \neq 0)
\]
Un exponente negativo \textbf{invierte} la base y cambia el exponente a positivo.
\end{tcolorbox}

\begin{tcolorbox}[ejemplobox, title={Ejemplos resueltos}]
\textbf{a)} $\left(\dfrac{2}{3}\right)^{-2} = \left(\dfrac{3}{2}\right)^{2} = \dfrac{9}{4}$\\[8pt]
\textbf{b)} $\left(\dfrac{-1}{4}\right)^{-3} = \left(\dfrac{4}{-1}\right)^{3} = (-4)^3 = -64$\\[8pt]
\textbf{c)} $3^{-2} = \dfrac{1}{3^2} = \dfrac{1}{9}$\\[8pt]
\textbf{d)} $\dfrac{3^{-1} + 4^{-1}}{5^{-1}} = \dfrac{\tfrac{1}{3} + \tfrac{1}{4}}{\tfrac{1}{5}} = \left(\dfrac{4+3}{12}\right) \cdot 5 = \dfrac{7}{12} \cdot 5 = \dfrac{35}{12}$
\end{tcolorbox}

\begin{tcolorbox}[ejerciciobox, title={Ejercicios — Potencias y Exponente Negativo}]
\begin{multicols}{2}
\begin{enumerate}[label=\textbf{\arabic*.}, itemsep=8pt]
  \item $\left(\dfrac{1}{2}\right)^3 =$
  \item $\left(\dfrac{-2}{3}\right)^2 =$
  \item $\left(\dfrac{3}{4}\right)^{-2} =$
  \item $5^{-3} =$
  \item $\left(\dfrac{-1}{2}\right)^{-4} =$
  \item $\left(\dfrac{2}{5}\right)^{-1} =$
  \item $\dfrac{2^{-1} + 3^{-1}}{6^{-1}} =$
  \item $\left(\dfrac{4}{3}\right)^{-2} \cdot \left(\dfrac{3}{2}\right)^2 =$
\end{enumerate}
\end{multicols}
\vspace{2pt}
%\small\textit{respuestas: 1) $\tfrac{1}{8}$ \; 2) $\tfrac{4}{9}$ \; 3) $\tfrac{16}{9}$ \; 4) $\tfrac{1}{125}$ \; 5) $16$ \; 6) $\tfrac{5}{2}$ \; 7) $5$ \; 8) $\tfrac{9}{16} \cdot \tfrac{9}{4}=\tfrac{81}{64}$}
\end{tcolorbox}

\subbloque{11. Operaciones combinadas con racionales}

\begin{tcolorbox}[ejemplobox, title={Ejemplo resuelto — nivel intermedio}]
Calcular: $\dfrac{3}{4} - \dfrac{2}{3} \cdot \dfrac{9}{8} + \dfrac{1}{2}$

\medskip
\textbf{Paso 1:} La multiplicación primero: $\dfrac{2}{3} \cdot \dfrac{9}{8} = \dfrac{18}{24} = \dfrac{3}{4}$

\textbf{Paso 2:} Luego la suma y resta: $\dfrac{3}{4} - \dfrac{3}{4} + \dfrac{1}{2} = 0 + \dfrac{1}{2} = \dfrac{1}{2}$
\end{tcolorbox}

\begin{tcolorbox}[ejemplobox, title={Ejemplo resuelto — con agrupaciones}]
Calcular: $\left(\dfrac{1}{2} - \dfrac{1}{3}\right) \div \left(\dfrac{1}{4} + \dfrac{1}{12}\right)$

\medskip
\textbf{Paso 1:} Paréntesis izquierdo: $\dfrac{1}{2} - \dfrac{1}{3} = \dfrac{3-2}{6} = \dfrac{1}{6}$

\textbf{Paso 2:} Paréntesis derecho: $\dfrac{1}{4} + \dfrac{1}{12} = \dfrac{3+1}{12} = \dfrac{4}{12} = \dfrac{1}{3}$

\textbf{Paso 3:} División: $\dfrac{1}{6} \div \dfrac{1}{3} = \dfrac{1}{6} \cdot \dfrac{3}{1} = \dfrac{3}{6} = \dfrac{1}{2}$
\end{tcolorbox}

\begin{tcolorbox}[ejerciciobox, title={Ejercicios — Nivel Intermedio (Operaciones Combinadas con Racionales)}]
\begin{enumerate}[label=\textbf{\arabic*.}, itemsep=8pt]
  \item $\dfrac{1}{2} + \dfrac{3}{4} \cdot \dfrac{2}{9} - \dfrac{1}{6} =$
  \item $\left(\dfrac{3}{5} - \dfrac{1}{10}\right) \cdot \dfrac{5}{4} + \dfrac{1}{8} =$
  \item $\dfrac{2}{3} \div \left(\dfrac{1}{2} + \dfrac{1}{6}\right) - \dfrac{1}{4} =$
  \item $\left(\dfrac{-2}{3}\right)^2 \cdot \dfrac{3}{2} - \dfrac{1}{3} =$
  \item $\dfrac{\tfrac{3}{4} - \tfrac{1}{2}}{\tfrac{1}{3} + \tfrac{1}{6}} =$
  \item $\dfrac{1}{2} \cdot \left[\dfrac{4}{3} - \left(\dfrac{1}{2}\right)^2\right] + \dfrac{5}{12} =$
\end{enumerate}
\vspace{2pt}
%\small\textit{respuestas: 1) $\tfrac{7}{12}$ \; 2) $\tfrac{3}{4}$ \; 3) $\tfrac{3}{4}$ \; 4) $\tfrac{1}{3}$ \; 5) $\tfrac{1}{2}$ \; 6) $1$}
\end{tcolorbox}

% ============================================================
% --- BLOQUE C: NIVEL AVANZADO RACIONALES ---
% ============================================================
\bloqueheader{C}{blockc}{Operaciones con Racionales — Nivel Avanzado}

\subbloque{12. Expresiones complejas con múltiples operaciones}

\begin{tcolorbox}[ejemplobox, title={Ejemplo resuelto — nivel avanzado}]
Calcular: $\left\{\dfrac{3}{4} - \left[\dfrac{1}{2} - \left(\dfrac{2}{3}\right)^2\right] \cdot 3\right\} \div \dfrac{-5}{8}$

\medskip
\textbf{Paso 1:} Paréntesis interno: $\left(\dfrac{2}{3}\right)^2 = \dfrac{4}{9}$

\textbf{Paso 2:} Corchete: $\dfrac{1}{2} - \dfrac{4}{9} = \dfrac{9-8}{18} = \dfrac{1}{18}$

\textbf{Paso 3:} Multiplicar por 3: $\dfrac{1}{18} \cdot 3 = \dfrac{3}{18} = \dfrac{1}{6}$

\textbf{Paso 4:} Llave: $\dfrac{3}{4} - \dfrac{1}{6} = \dfrac{9-2}{12} = \dfrac{7}{12}$

\textbf{Paso 5:} División: $\dfrac{7}{12} \div \dfrac{-5}{8} = \dfrac{7}{12} \cdot \dfrac{8}{-5} = \dfrac{56}{-60} = \dfrac{-14}{15}$
\end{tcolorbox}

\begin{tcolorbox}[ejerciciobox, title={Ejercicios — Nivel Avanzado (Racionales con Múltiples Agrupaciones)}]
Resuelve completamente, mostrando cada paso:
\begin{enumerate}[label=\textbf{\arabic*.}, itemsep=12pt]
  \item $\left\{\dfrac{5}{6} - \left[\dfrac{2}{3} - \left(\dfrac{1}{4} + \dfrac{1}{3}\right)\right]\right\} \cdot \left(-\dfrac{6}{5}\right) =$

  \item $\dfrac{\left(\dfrac{3}{4}\right)^{-1} - \left(\dfrac{2}{3}\right)^{-1}}{\dfrac{1}{2} - \dfrac{1}{3}} =$

  \item $\left(\dfrac{-1}{2}\right)^3 + \left[\dfrac{3}{4} - \dfrac{1}{2} \cdot \left(\dfrac{2}{3} - 1\right)\right] - \dfrac{5}{12} =$

  \item $\left\{\left[\dfrac{-2}{3} + \dfrac{5}{6}\right]^2 - \dfrac{1}{4}\right\} \div \left[\dfrac{3}{8} - \left(\dfrac{-1}{4}\right)\right] =$

  \item $\dfrac{3}{4} \cdot \left\{\dfrac{-8}{3} + \left[\dfrac{1}{2} - \dfrac{1}{6}\right]^{-2} - 2\right\} =$

  \item $\dfrac{\left(\dfrac{5}{3} - \dfrac{3}{2}\right)^2}{\left(\dfrac{1}{6}\right)^{-1} - 5} + \dfrac{2}{3} \cdot \dfrac{-9}{4} =$
\end{enumerate}
\vspace{4pt}
\begin{tcolorbox}[colback=white, colframe=zenitnavy, arc=2pt, left=4pt, right=4pt, top=3pt, bottom=3pt]
\textbf{Solución detallada del \#2:}
\[
\dfrac{\left(\dfrac{3}{4}\right)^{-1} - \left(\dfrac{2}{3}\right)^{-1}}{\dfrac{1}{2} - \dfrac{1}{3}}
= \dfrac{\dfrac{4}{3} - \dfrac{3}{2}}{\dfrac{1}{6}}
= \dfrac{\dfrac{8-9}{6}}{\dfrac{1}{6}}
= \dfrac{\dfrac{-1}{6}}{\dfrac{1}{6}}
= \dfrac{-1}{6} \cdot \dfrac{6}{1}
= \mathbf{-1}
\]
\end{tcolorbox}
\end{tcolorbox}

\newpage

% ============================================================
% --- EJERCICIOS FINALES: MEZCLA DE TODOS LOS TEMAS ---
% ============================================================
\bloqueheader{D}{blockc}{Ejercicios Finales — Mezcla de Todos los Temas}

\begin{tcolorbox}[instrbox]
Los siguientes ejercicios mezclan \textbf{todos los contenidos}: enteros, fracciones, potencias, jerarquía y agrupaciones. Se incluyen ejercicios tipo PSU/SIMCE para máxima preparación.
\end{tcolorbox}

\vspace{6pt}

\begin{enumerate}[label=\textbf{\arabic*.}, leftmargin=*, itemsep=14pt]

\item Calcula:
\[
  (-3)^2 + 2 \cdot (-4) - \dfrac{(-12)}{3} =
\]
\begin{enumerate}[label=\Alph*)]
  \item $-7$ \quad B) $1$ \quad C) $5$ \quad D) $13$ \quad E) $17$
\end{enumerate}

\item ¿Cuánto vale la siguiente expresión?
\[
  \dfrac{2^{-1} + \left(\dfrac{1}{3}\right)^{-1}}{(-2)^3 + 10}
\]
\begin{enumerate}[label=\Alph*)]
  \item $\dfrac{1}{7}$ \quad B) $\dfrac{7}{2}$ \quad C) $\dfrac{7}{4}$ \quad D) $\dfrac{7}{6}$ \quad E) $3$
\end{enumerate}

\item Simplifica y calcula:
\[
  \dfrac{3}{8} \div \dfrac{9}{4} + \left(\dfrac{-1}{2}\right)^3 \cdot \dfrac{4}{3} =
\]

\item ¿Cuál es el resultado de?
\[
  \left\{(-2)^4 - [3 \cdot (-1 + 5)^2 - 20]\right\} \cdot \dfrac{-3}{4}
\]

\item Calcula el valor de la expresión, sabiendo que es un número entero:
\[
  \left\{\left[\dfrac{5}{4} - \left(\dfrac{-3}{2}\right)^2\right] \cdot (-8) + \dfrac{1}{3}\right\} \cdot 6 =
\]

\item Determina cuál de las siguientes expresiones tiene el mayor valor:
\begin{enumerate}[label=\Alph*)]
  \item $(-2)^6 \div (-4)$
  \item $\dfrac{(-3)^3 + 5}{(-2)^2 - 6}$
  \item $\left(\dfrac{-3}{4}\right)^{-2}$
  \item $3^{-1} - (-3)^{-1}$
  \item $\dfrac{(-1)^{100} + (-1)^{101}}{(-1)^{99}}$
\end{enumerate}

\item Calcula:
\[
  \dfrac{-5}{6} - \dfrac{3}{4} \cdot \left(\dfrac{-2}{3}\right)^2 + \dfrac{1}{3} \div \dfrac{-4}{9}
\]

\item Resuelve la siguiente expresión mostrando todos los pasos:
\[
  \left\{(-3)^2 \cdot \left[\dfrac{-2}{3} + \dfrac{5}{6}\right] - \left(\dfrac{1}{2}\right)^{-2}\right\} \div \left(-\dfrac{5}{4}\right)
\]

\item Si $a = -\dfrac{2}{3}$ y $b = \dfrac{3}{4}$, calcula $a^2 - 2ab + b^2$.

\item Indica si la siguiente igualdad es verdadera o falsa y justifica:
\[
  \dfrac{(-2)^3 + (-3)^2}{(-1)^5 \cdot 2^2} = \dfrac{1}{4}
\]

\item Calcula:
\[
  \dfrac{\left(\dfrac{3}{2}\right)^{-2} - \left(\dfrac{2}{3}\right)^{2}}{\left(\dfrac{3}{2}\right)^{2} + \left(\dfrac{2}{3}\right)^{-2}}
\]

\item \textbf{(Desafío)} Calcula:
\[
  \left\{-\left[(-2)^3 - \dfrac{3}{4} \cdot \left(\dfrac{-4}{3}\right)^2 + \left(\dfrac{1}{2}\right)^{-3}\right]^2 + 5 \cdot (-1)^{101}\right\} \div \dfrac{-3}{2}
\]

\end{enumerate}

\vspace{10pt}

\begin{tcolorbox}[instrbox]
\textbf{Respuestas Ejercicios Finales:}
\begin{enumerate}[leftmargin=*, itemsep=3pt]
  \item[\textbf{1.}] $9 + (-8) - (-4) = 9 - 8 + 4 = 5$ \quad $\Rightarrow$ \textbf{C)}
  \item[\textbf{2.}] $\dfrac{\tfrac{1}{2}+3}{-8+10} = \dfrac{\tfrac{7}{2}}{2} = \dfrac{7}{4}$ \quad $\Rightarrow$ \textbf{C)}
  \item[\textbf{3.}] $\dfrac{3}{8} \cdot \dfrac{4}{9} + \dfrac{-1}{8} \cdot \dfrac{4}{3} = \dfrac{12}{72} - \dfrac{4}{24} = \dfrac{1}{6} - \dfrac{1}{6} = \mathbf{0}$
  \item[\textbf{4.}] $\{16 - [3 \cdot 16 - 20]\} \cdot \dfrac{-3}{4} = \{16 - 28\} \cdot \dfrac{-3}{4} = (-12)\cdot\dfrac{-3}{4} = \mathbf{9}$
  \item[\textbf{5.}] $\left\{\left[\dfrac{5}{4} - \dfrac{9}{4}\right] \cdot (-8) + \dfrac{1}{3}\right\} \cdot 6 = \{(-1) \cdot (-8) + \tfrac{1}{3}\} \cdot 6 = \{8 + \tfrac{1}{3}\} \cdot 6 = \tfrac{25}{3} \cdot 6 = \mathbf{50}$
  \item[\textbf{7.}] $\dfrac{-5}{6} - \dfrac{3}{4} \cdot \dfrac{4}{9} + \dfrac{1}{3} \cdot \dfrac{-9}{4} = \dfrac{-5}{6} - \dfrac{1}{3} - \dfrac{3}{4} = \dfrac{-10-4-9}{12} = \dfrac{-23}{12}$
  \item[\textbf{9.}] $a^2 - 2ab + b^2 = (a-b)^2 = \left(-\dfrac{2}{3} - \dfrac{3}{4}\right)^2 = \left(\dfrac{-8-9}{12}\right)^2 = \left(\dfrac{-17}{12}\right)^2 = \dfrac{289}{144}$
  \item[\textbf{10.}] V.F.: $\dfrac{-8+9}{(-1)\cdot 4} = \dfrac{1}{-4} = -\dfrac{1}{4} \neq \dfrac{1}{4}$ \quad $\Rightarrow$ \textbf{FALSA}
\end{enumerate}
\end{tcolorbox}

\vspace{8pt}
\begin{center}
  \rule{\linewidth}{1.5pt}\\[4pt]
  {\large\bfseries\color{zenitnavy} ¡Bien hecho! La constancia y la práctica son la clave del éxito.}\\[2pt]
  {\color{zenitgold}\small Academia Zenit · Matemática 1° Medio}
  \rule{\linewidth}{1.5pt}
\end{center}

\end{document}
